% ============================================================
% CV — Atharva Vijay Shinde
% Research / Generative AI / Data Engineering
% Compile with: lualatex main_example.tex
% ============================================================

\documentclass[11pt,a4paper,sans]{moderncv}

% ------------------------------------------------------------
% STYLE
% ------------------------------------------------------------
\moderncvstyle{banking}
\moderncvcolor{blue}

\renewcommand*{\namefont}{\fontsize{34}{36}\bfseries\upshape}
\colorlet{firstnamecolor}{color1}
\colorlet{lastnamecolor}{color1}

\renewcommand*{\sectionstyle}[1]{{\sectionfont\color{color1}#1}}

% ------------------------------------------------------------
% PACKAGES
% ------------------------------------------------------------
\usepackage[utf8]{inputenc}
\usepackage[scale=0.78]{geometry}

\AtEndPreamble{
    \hypersetup{
        colorlinks=true,
        linkcolor=blue,
        filecolor=magenta,
        urlcolor=blue,
        pdftitle={Atharva Vijay Shinde - CV},
        pdfauthor={Atharva Vijay Shinde},
        pdfpagemode=UseNone
    }
}

% ------------------------------------------------------------
% PERSONAL INFORMATION
% ------------------------------------------------------------
\name{Atharva}{Shinde}

\address{Nagpur}{India}{}

\phone[mobile]{+91-8308257250}

\email{sbjitatharvas@gmail.com}

\social[linkedin]{atharva-shinde7474}

\extrainfo{
    \href{https://github.com/Atharva680}{GitHub}
    \textbar\ 
    \href{https://atharva-portfolio-betanine.vercel.app}{Portfolio}
}

% ============================================================
% DOCUMENT
% ============================================================
\begin{document}

\makecvtitle

% ============================================================
% RESEARCH STATEMENT
% ============================================================

\section{Research Statement}

Final-year B.E. Computer Science Engineering student graduating July 2026
with a strong interest in Artificial Intelligence, Large Language Models,
software engineering, and applied research. Experienced in building
production-grade LLM systems, enterprise RAG pipelines, AI agents, and
large-scale data pipelines. Interested in the intersection of LLMs,
programming languages, formal software engineering methods, code analysis,
and automated code transformation.

% ============================================================
% RESEARCH ALIGNMENT
% ============================================================

\section{Research Alignment -- FoSSE Focus Areas}

\begin{itemize}

\item
\textbf{AI for Code Analysis \& Transformation:}
Built AI agent systems performing automated document/code-like migration,
classification, extraction, and structured transformation tasks using LLMs.
Experienced in prompting LLMs for structured code-like output and analysis.

\item
\textbf{LLM-Powered Software Engineering Tools:}
Designed and deployed enterprise-scale LLM applications using LangChain
and OpenAI API that process 100,000+ structured documents with 95\%
retrieval accuracy, demonstrating the ability to build reliable AI tools
for large-scale repositories.

\item
\textbf{Static Analysis \& Data Pipeline Research:}
Engineered end-to-end data pipelines for automated ingestion, parsing,
chunking, embedding, indexing, and semantic analysis of large document
collections. These skills are transferable to static analysis of large
codebases.

\item
\textbf{Python \& ML for Research Experimentation:}
Developed modular Python research code with MLflow experiment tracking,
including parameter, metric, artifact, and model-version tracking for
reproducible experimentation.

\item
\textbf{Formal Methods Foundation:}
Academic background in algorithms, data structures, discrete mathematics,
theory of computation, database systems, operating systems, and software
engineering, providing foundations for programming-language and formal
verification research.

\end{itemize}

% ============================================================
% TECHNICAL SKILLS
% ============================================================

\section{Technical Skills}

\cvitem{Programming}
{Python, SQL, JavaScript, Shell}

\cvitem{AI \& LLM Systems}
{LangChain, OpenAI API, RAG Pipelines, Agentic AI, Prompt Engineering,
Embeddings, Vector Search}

\cvitem{Machine Learning}
{PyTorch, TensorFlow, Scikit-learn, LSTM, Deep Q-Learning,
NLP, Computer Vision, MLflow}

\cvitem{Software Engineering}
{FastAPI, REST APIs, Docker, Microservices, GitHub, CI/CD,
Agile, Code Reviews, OOP, Testing}

\cvitem{Data Engineering \& Cloud}
{Apache Spark, Databricks, Delta Lake, ETL/ELT,
Microsoft Azure, Batch Processing, Streaming}

\cvitem{Research Tools}
{MLflow, PostgreSQL, PGVector, Pinecone, Jupyter, Git}

\cvitem{Platforms}
{Linux (Ubuntu), Windows, Microsoft Azure, GitHub, VS Code}

% ============================================================
% PROFESSIONAL EXPERIENCE
% ============================================================

\section{Professional Experience}

\cventry
{Nov 2025 -- Apr 2026}
{Generative AI \& Data Engineering Intern}
{InfoCepts -- Data \& AI}
{Nagpur, India}
{}
{
\begin{itemize}

\item
Built enterprise RAG pipelines using LangChain and OpenAI API that
automatically ingest, parse, chunk, embed, and semantically search
100,000+ enterprise documents.

\item
Designed multi-step Agentic AI systems with dynamic tool-calling to
autonomously perform classification, extraction, and structured
transformation tasks.

\item
Optimized retrieval and vector-search strategies using PGVector,
reducing OpenAI API costs by 25\%.

\item
Packaged AI solutions into production-grade FastAPI microservices
running in Docker containers on Microsoft Azure.

\item
Engineered end-to-end ETL/ELT pipelines using Python, Apache Spark,
and Delta Lake on Databricks for large-scale structured data processing.

\item
Tracked ML experiments using MLflow, recording parameters, metrics,
artifacts, and model versions to support reproducibility.

\item
Shipped a real-time face recognition pipeline using OpenCV, YuNet,
and TensorFlow achieving 98\% accuracy and sub-100ms latency.

\item
Received a Certificate of Recognition at InfoCepts GALA 2026 for
outstanding applied AI engineering work.

\end{itemize}
}

\cventry
{Jun 2024 -- Aug 2024}
{Web Developer Intern}
{LGPS Hybrid Energy Pvt. Ltd.}
{Nagpur, India}
{}
{
\begin{itemize}

\item
Built a full-stack energy monitoring dashboard using Node.js,
React, and Chart.js with live REST API integration.

\item
Developed practical full-stack software engineering solutions
for real-time energy monitoring and visualization.

\end{itemize}
}

\cventry
{2023 -- 2025}
{Placement Coordinator}
{S.B. Jain Institute of Technology, Management \& Research}
{Nagpur, India}
{}
{
\begin{itemize}

\item
Coordinated 15+ company recruitment drives involving organizations
including HCL and Amazon.

\item
Managed recruitment logistics, communication, coordination, and
student-company interactions.

\end{itemize}
}

% ============================================================
% RESEARCH-RELEVANT PROJECTS
% ============================================================

\section{Research-Relevant Projects}

\cventry
{}
{Enterprise RAG System over Large Document Repository}
{LangChain | OpenAI API | PGVector | FastAPI | Python | Azure}
{}
{}
{
\begin{itemize}

\item
Production-oriented system performing automated semantic analysis
over 100,000+ documents through chunking, embedding, vector indexing,
retrieval, and LLM-powered response generation.

\item
Achieved 95\% retrieval accuracy and reduced API costs by 25\%
through optimized retrieval strategies.

\end{itemize}
}

\cventry
{}
{Agentic AI Code-Like Transformation Platform}
{LangChain | Python | FastAPI | Docker | Azure}
{}
{}
{
\begin{itemize}

\item
Developed a multi-step AI agent capable of autonomous structured
document classification, extraction, and transformation using LLMs.

\item
Automated complex structured analysis workflows, saving 260+ hours
annually.

\end{itemize}
}

\cventry
{}
{LSTM + Reinforcement Learning Research Experiment}
{PyTorch | TensorFlow | Deep Q-Learning | MLflow}
{}
{}
{
\begin{itemize}

\item
Developed an LSTM forecasting model achieving RMSE of 0.0257 kWh
and integrated a Deep Q-Learning agent.

\item
Tracked experiments using MLflow, covering experiment parameters,
metrics, artifacts, and model versions for reproducibility.

\end{itemize}
}

% ============================================================
% EDUCATION
% ============================================================

\section{Education}

\cventry
{Jul 2022 -- Jul 2026}
{Bachelor of Engineering (B.E.) -- Computer Science Engineering}
{S.B. Jain Institute of Technology, Management and Research (SBJITMR)}
{Nagpur, India}
{}
{
\textbf{RTM Nagpur University}

\textbf{CGPA: 7.54/10.0 \quad | \quad Grade: A+}

\textbf{Relevant Coursework:}
Data Structures \& Algorithms, Theory of Computation,
Discrete Mathematics, Database Systems, Machine Learning,
Artificial Intelligence, Operating Systems, Computer Networks,
Software Engineering
}

% ============================================================
% CERTIFICATIONS
% ============================================================

\section{Certifications}

\begin{itemize}

\item
\textbf{Databricks Certified Generative AI Engineer Associate}
-- Credential: 181091666 -- Valid Apr 2026 -- Apr 2028

\item
\textbf{Databricks Certified Data Engineer Associate}
-- Credential: 180763435 -- Valid Apr 2026 -- Apr 2028

\item
\textbf{Databricks Generative AI Fundamentals}
-- April 2026

\item
\textbf{Cyber Security \& Ethical Hacking}
-- National Skill Development Corporation -- July 2025

\item
\textbf{IQS Python Programming}
-- Grade: A+

\end{itemize}

% ============================================================
% ACHIEVEMENTS & ACADEMIC ACTIVITIES
% ============================================================

\section{Achievements \& Academic Activities}

\begin{itemize}

\item
\textbf{Certificate of Recognition -- InfoCepts GALA 2026:}
Recognised for outstanding applied AI engineering during production
internship at a global Data \& AI company.

\item
\textbf{Department Award Trophy -- IEEE Computer Society Chapter (2025):}
Led 10+ technical AI \& Robotics events, demonstrating technical
leadership, initiative, and research communication.

\item
\textbf{Campus Ambassador -- IIT Guwahati (May -- Sep 2025):}
Engaged with researchers and students at a premier national research
institution and participated in academic and research community
networking.

\item
\textbf{Placement Coordinator -- SBJITMR (2023 -- 2025):}
Coordinated 15+ company recruitment drives involving HCL, Amazon,
and other organizations.

\end{itemize}

% ============================================================
% LANGUAGES
% ============================================================

\section{Languages}

\cvitem{}{English (Professional) \quad | \quad Hindi (Native)
\quad | \quad Marathi (Native)}

% ============================================================
% END
% ============================================================

\end{document}